\chapter{Literature Review}
\label{chap:literature}

This chapter surveys the prior art the thesis builds on, grouped into
seven themes. The first two — classical truss optimisation and
evolutionary metaheuristics — are the structural-engineering and
optimisation-theory foundations on which every modern truss study
still stands. The next three — neural surrogates, reinforcement
learning, and LLM-assisted design — are the machine-learning
accelerators this thesis combines. \Cref{sec:lit:is800} documents
the Indian design code (IS~800:2007) that anchors the compliance
layer, and \cref{sec:lit:gap} synthesises the literature into a
one-paragraph statement of the research gap this thesis fills.

\section{Classical Truss Optimisation}
\label{sec:lit:classical}

Truss optimisation as a formal mathematical problem dates to the
early 1960s, when Louis Schmit argued that structural design should
be posed as a constrained-nonlinear-programming
problem~\cite{schmit1974}. Schmit and Farshi's 1974 paper is
canonically credited with introducing the modern formulation: a
vector of cross-sectional areas $\mathbf{x}$ minimising
total weight $W(\mathbf{x}) = \sum_i \rho A_i L_i$ subject to
per-member stress inequalities $|\sigma_i(\mathbf{x})| \le
\sigma_{\text{allow}}$ and per-node displacement inequalities
$|u_j(\mathbf{x})| \le u_{\text{allow}}$. Both the objective and
constraints depend on the areas through a linear elastic
finite-element analysis. That same paper also introduced the 10-bar
planar truss that has since become the single most-cited benchmark
in the field; the published optimum of \num{5060.85}~lb has been
matched to four significant figures by dozens of subsequent studies.

Venkayya's 1971 optimality-criteria method~\cite{venkayya1971}
produced the next canonical benchmark — the 25-bar spatial truss —
and established that gradient-based algorithms could match
hand-computed optima on small- to medium-scale sizing problems.
Fleury and Schmit in 1980~\cite{fleury1980} extended the approach
to the 72-bar four-storey spatial tower, which exposed the
sensitivity of optimisation outcomes to group-symmetry choices and
two-load-case interactions that simpler benchmarks had hidden.
Schmit and Miura's 1976 stress-ratio method~\cite{schmit1976} gave
a dual-method framework for tension-dominated trusses, and Sunar
and Belegundu's 1991 paper~\cite{sunar1991} revisited the 10-bar
benchmark with improved numerics, producing the \num{5060.85}~lb
value used throughout this thesis.

Erbatur et al.~\cite{erbatur2000} in 2000 reported a genetic-
algorithm comparison across the 10-, 25-, and 72-bar benchmarks and
published discrete-area optima that are still standard reference
points. Camp and Bichon~\cite{camp2004} in 2004 applied ant-colony
optimisation to the 72-bar tower and reported the
\num{380.24}~lb figure that every subsequent metaheuristic paper
has either tried to beat or match. Kaveh and Talatahari's 2010
large-scale study~\cite{kaveh2010} established the 200-bar planar
truss as the stress-scalability benchmark. Bekdaş et
al.~\cite{bekdas2015} applied flower-pollination to the 72-bar
tower and reported a continuous optimum of \num{379.62}~lb that
remains the most-cited continuous target.

Two observations from this history are relevant to the present
thesis. First, the benchmark set has barely changed in fifty years;
any new method is tested against the same 10-, 25-, 72-, and
200-bar problems, which makes cross-paper comparison possible.
Second, published ``optima'' for the 72-bar tower are sensitive to
the constraint-handling convention. Phase~7.5 of this thesis (see
\cref{sec:results:72bar}) shows that the published Camp-Bichon and
Bekdaş optima are infeasible by $22\,\%$ on stress and $62\,\%$ on
lateral-tip displacement when the constraint is enforced rigorously
in a modern FEM. The literature minima are reachable only with
soft-penalty constraint handling that admits small violations. This
distinction has methodological consequences that the thesis makes
explicit in \cref{chap:results}.

\section{Evolutionary and Swarm Metaheuristics}
\label{sec:lit:metaheuristics}

Evolutionary computation as applied to engineering design was first
argued for as an alternative to gradient methods on non-convex,
non-differentiable, mixed-variable problems where classical
calculus-based methods either fail or require heavy smoothing.
Deb's 2001 textbook~\cite{deb2001} remains the standard reference
for the single- and multi-objective genetic-algorithm machinery
used in this thesis. The Simulated Binary Crossover (SBX) and
polynomial mutation operators, with the distribution indices
$\eta_c = 15$ and $\eta_m = 20$ used here, originate in the same
text.

Kennedy and Eberhart~\cite{kennedy1995pso} introduced Particle
Swarm Optimisation in 1995 as a population-based search inspired by
bird-flocking dynamics. PSO has historically been competitive with
GA on continuous-variable sizing problems and is often the method
of choice when gradient information is unavailable but the
objective is smooth. This thesis uses PSO with the default pymoo
hyperparameters and finds it notably tighter than GA in cross-seed
variance on the 25-bar and 72-bar benchmarks (standard deviation
below $0.1$~lb, versus GA's $0.7$~lb on the 72-bar).

Multi-objective extension is provided by Deb, Pratap, Agarwal and
Meyarivan's NSGA-II~\cite{deb2002nsga2}, published in 2002 and the
most-cited evolutionary-algorithm paper in structural optimisation.
NSGA-II's non-dominated sorting plus crowding-distance selection
has become the standard against which every subsequent multi-
objective evolutionary algorithm is measured. For the weight-
versus-displacement Pareto front the thesis reports in
\cref{sec:results:pareto}, NSGA-II was chosen without much
deliberation — it is the field-standard baseline.

Constraint handling in evolutionary algorithms is a sub-field in
its own right. Coello Coello's 2002 survey~\cite{coello2002}
catalogues the approaches: death penalty, static penalty, dynamic
penalty, adaptive penalty, feasibility rules, and multi-objective
constraint handling. This thesis uses the feasibility-rule approach
implicit in pymoo's constraint handling (Blank and
Deb~\cite{blank2020pymoo}): a feasible solution is always preferred
over an infeasible one; between two infeasible solutions the less-
violating one is preferred; between two feasible solutions the
better-objective one wins. The constraint expression
$g(\mathbf{x}) \le 0$ is chosen so that the value of $g$ at any
infeasible point is a meaningful violation magnitude.

Blank and Deb~\cite{blank2020pymoo} published pymoo in 2020 as a
unified Python framework for multi-objective optimisation. Every
optimiser in this thesis is a thin wrapper over pymoo's
implementation. This was a deliberate choice: re-implementing GA,
PSO, or NSGA-II from scratch would add thousands of lines of
algorithmic code without contributing anything the literature
needs. The thesis's integrative novelty is in the layer above
pymoo, not at the algorithm-engine level.

\section{Neural Surrogate Models for Structural Analysis}
\label{sec:lit:surrogate}

The use of machine-learning models as cheap approximations to
expensive physics-based simulations predates deep learning. Queipo
et al.~\cite{queipo2005} synthesised the surrogate-based analysis
and optimisation literature in 2005, cataloguing polynomial
response-surfaces, radial-basis functions, Kriging, and support-
vector regression as the then-dominant approximators. Forrester,
Sóbester, and Keane's 2008 textbook~\cite{forrester2008}
established Kriging (Gaussian process regression) as the surrogate
of choice for smooth low-dimensional problems and gave the
infill-criterion framework (expected improvement) that turned
surrogates from one-shot approximators into sequential design-of-
experiments drivers.

The post-2015 deep-learning wave shifted surrogate modelling toward
multi-layer perceptrons (MLPs) and, more recently, graph neural
networks. Goodfellow, Bengio, and Courville's
textbook~\cite{goodfellow2016dl} is the standard reference for the
training machinery used here — Adam's stochastic-gradient
optimiser~\cite{kingma2015adam}, ReLU activations, dropout
regularisation, early stopping on held-out validation loss.

For structural engineering specifically, Sun, Wang, and Wu's 2023
work~\cite{sun2023} combined a deep neural network with a genetic
algorithm for composite-structure optimisation and reported
speedups consistent with the $50$--$150\times$ range this thesis
targets. Persia et al.~\cite{persia2025} trained neural surrogates
for topology-optimisation post-processing, a related but distinct
problem.

The architecture this thesis uses — a three-hidden-layer MLP of
widths $[256, 128, 64]$ with multi-head output (weight,
displacement, stress, feasibility logit) — is deliberately
unglamorous. Truss-sizing problems have at most $\sim 30$ design
variables (200-bar benchmark) and smooth outputs, so there is no
prior reason to expect graph neural networks or transformers to
outperform a well-regularised MLP. The design choice is justified
post-hoc in \cref{sec:results:surrogate} by the achieved
$R^2 > 0.98$ on held-out weight prediction.

\section{Reinforcement Learning in Structural Design}
\label{sec:lit:rl}

Reinforcement learning applied to structural design is a recent
development. Sutton and Barto's 2018 textbook~\cite{sutton2018rl}
is the canonical theoretical reference; Schulman et al.'s 2017
Proximal Policy Optimisation paper~\cite{schulman2017ppo} is the
practical workhorse this thesis uses. PPO's clipped-surrogate
objective is stable enough to train out-of-the-box on sparse-reward
continuous-action problems, which is what truss design becomes
once it is framed as an episodic task.

Hayashi and Ohsaki's 2020 study~\cite{hayashi2020rl} is the direct
structural-design precedent: they trained a deep Q-network to
select members from a discrete cross-section catalogue for planar
truss sizing. Brown and Mueller~\cite{brown2022rl} applied PPO to
topology optimisation of 2D continuum structures in 2022.
Rochefort-Beaudoin et al.'s 2024 open-source SOgym
environment~\cite{rochefort2024sogym} provided the first
standardised Gymnasium-compatible RL benchmark for structural
optimisation. Zhao et al.~\cite{zhao2021rl} applied RL to
topology-optimisation problems and established the coupling of an
RL agent with a surrogate solver (our Layer~3) as a tractable
training strategy.

The thesis's RL layer uses PPO as the agent and the Layer~3
surrogate as the environment's physics. Training RL directly
against the FEM is computationally infeasible on a laptop —
$\num{100000}$ timesteps $\times$ \num{50}~ms per FEM solve would
exceed \num{80} minutes of wall-clock per benchmark. With the
surrogate substituted for the FEM, training completes in under
five minutes. This surrogate-in-environment pattern is, as far as
the author can find, the closest structural-engineering analogue
to the sim-to-real RL work common in robotics.

\section{LLM-Assisted Engineering Design}
\label{sec:lit:llm}

Large language models have entered the structural-design literature
in the last two years. Two precedents directly motivated the
Layer~5 (LLM warm-start) contribution of this thesis. Geng et
al.'s 2024 paper~\cite{geng2024} combined an LLM with the
OpenSeesPy structural-analysis package to generate analysis scripts
from natural-language descriptions of loading scenarios; the LLM
acts as a pre-processor that translates intent into an analysis-
ready model. Makatura et al.~\cite{makatura2024} surveyed LLMs for
design and manufacturing broadly, arguing that LLMs should be
positioned as semantic co-pilots that bridge between human intent
and formal optimisation specifications, rather than as black-box
design generators that replace the optimisation step.

This thesis follows Makatura et al.'s framing. The LLM (Anthropic
Claude~\cite{anthropic-docs}) is asked for a single initial design
vector that warm-starts the GA; it does not replace the GA, and its
output is never reported as a final answer. Wei et al.'s 2022
chain-of-thought prompting paper~\cite{wei2022cot} motivates the
prompt structure, which asks the LLM to narrate its reasoning
before emitting the final JSON. The reasoning text is archived
with every cached response, both for reproducibility and because
it occasionally surfaces a structural insight — on the 10-bar
benchmark the LLM explicitly identifies the redundant-member set
$\{2, 5, 6, 10\}$ and sizes them at the lower bound, matching what
an experienced structural engineer would do.

The empirical finding of \cref{sec:results:llm} — that LLM warm-
start helps most on small or structurally-sparse benchmarks and
provides no measurable benefit on the larger 72-bar problem — is,
as far as the author is aware, not previously reported. It
provides a scaling law of sorts for LLM-assisted design: the warm-
start's value grows with the complexity of the structural insight
the LLM can offer, not with the raw size of the optimisation
problem.

\section{Indian Standard IS~800:2007 Design Code}
\label{sec:lit:is800}

The Bureau of Indian Standards' IS~800:2007~\cite{is800-2007} is
the primary structural-steel design standard in use in India. It
superseded the working-stress-method-based IS~800:1984 by
introducing the limit-state-design (LSD) framework, aligning
Indian practice with the European and American families of codes
(Eurocode~3 and AISC~360 respectively). The thesis implements
three groups of clauses directly:
\begin{itemize}
  \item \textbf{Clause~3.8} — slenderness limits
        ($\lambda \le \num{180}$ for compression,
        $\lambda \le \num{400}$ for tension);
  \item \textbf{Clause~5.6.1} — deflection limits (typically span
        divided by \num{300} for roof trusses under live load);
  \item \textbf{Clauses~6.2 and~6.3} — tension design strength
        (gross-section yield and net-section rupture);
  \item \textbf{Clause~7.1} — compression design strength using the
        four column-buckling curves ($a$, $b$, $c$, $d$) derived
        from Perry-Robertson.
\end{itemize}
The standard steel-section catalogue SP~6(1):1964~\cite{sp6-1-1964}
provides rolled-section tables (ISA angles, ISMB beams, ISMC
channels) that the discrete-variable version of the benchmarks
draws from. Subramanian's 2008 textbook~\cite{subramanian2008} and
Duggal's 2014 textbook~\cite{duggal2014} provide worked examples
that were used as cross-checks against the thesis's implementation
of each clause.

The compliance layer (\texttt{src/constraints/is800\_checks.py})
treats each clause as a pure function taking member geometry,
stress, and load vectors and returning a boolean pass/fail. A
composite \texttt{full\_is800\_check} aggregates these into the
constraint vector $g(\mathbf{x})$ that the optimiser consumes.

\section{Research Gap and Thesis Positioning}
\label{sec:lit:gap}

Classical truss-sizing optimisation is mature; evolutionary and
swarm metaheuristics are mature; neural surrogates for structural
analysis have been demonstrated in multiple independent studies;
reinforcement learning for structural design has one small
precedent (Hayashi and Ohsaki, 2020); LLM-assisted structural
engineering has two precedents (Geng 2024, Makatura 2024). What
does not exist in the published literature is a single,
reproducible framework that combines all of these in one codebase,
validated against the standard four-benchmark set, and wired
throughout to the Indian IS~800:2007 design code. This thesis fills
that gap. The contribution is integrative rather than algorithmic:
the optimisers, the neural architecture, and the RL algorithm are
all standard choices; the novelty lies in the composition, in the
rigorous feasibility discipline applied to the 72-bar benchmark,
and in the first published scaling observation for LLM warm-start
on truss problems of increasing complexity.
